\chapter{Scenario Analysis}
\label{ch:scenarios}

\section{Concept}

The scenario engine allows comparing different ``what-if'' situations against a baseline. Each scenario defines a set of parameter modifications that are applied before re-running the solver. All scenarios share the same system structure --- only parameter values change.

\section{Configuration (Sheet 5\_1)}

Scenarios are defined in the \texttt{5\_1\_Scenario\_Definitions} sheet:

\begin{table}[h]
\centering
\small
\begin{tabular}{lllllll}
\toprule
\textbf{Scenario} & \textbf{Parameter} & \textbf{Type} & \textbf{Value} & \textbf{Start} & \textbf{End} \\
\midrule
Scenario\_A & TC\_05\_06 & multiply & 1.2 & 2030 & 2050 \\
Scenario\_A & F\_00\_01 & add & 200 & 2025 & 2040 \\
Scenario\_B & TC\_05\_06 & set & 0.9 & 2030 & 2050 \\
\bottomrule
\end{tabular}
\caption{Scenario definition example.}
\end{table}

\subsection{Modification Types}

\begin{description}
  \item[set] Replace the parameter with a fixed value
  \item[multiply] Multiply the current value by a factor
  \item[add] Add a value to the current parameter
\end{description}

\subsection{Time Range}

Each modification specifies a \texttt{Start\_Year} and \texttt{End\_Year}. The modification is only applied within this range. Outside the range, the baseline value is used.

\section{Supported Parameters}

The scenario engine auto-detects the parameter type from the naming pattern:

\begin{itemize}
  \item \textbf{Transfer coefficients} --- Names starting with \texttt{TC\_}
  \item \textbf{Flow data} --- Names starting with \texttt{F\_}
  \item \textbf{DSM parameters} --- Names containing \texttt{DSM\_}
  \item \textbf{FOMP parameters} --- Names containing \texttt{FOMP\_}
\end{itemize}

\section{Running Scenarios}

\begin{lstlisting}[style=python]
all_scenario_results = run_all_scenarios(
    mfa_system_configured=mfa_system,
    scenario_definitions=scenario_definitions,
    process_logic_map=process_logic_map,
    flow_tc_map=flow_tc_map,
    dsm_params_config=dsm_params,
    fomp_params_config=fomp_params,
)
\end{lstlisting}

Each scenario creates a deep copy of the baseline system, applies modifications, and re-runs the solver. Results are stored per scenario for comparison.

\section{Comparison Plots}

The scenario comparison visualization provides:
\begin{itemize}
  \item \textbf{Flow dynamics} --- Baseline vs.\ scenario flows over time
  \item \textbf{Stock dynamics} --- Baseline vs.\ scenario stocks over time
  \item Multi-scenario selection (compare 2+ scenarios simultaneously)
  \item Element selector (view any tracked element)
\end{itemize}

\begin{tipbox}
Use scenarios to evaluate policy options (e.g., increased recycling rates), technology changes (e.g., new processing efficiency), or demand projections (e.g., reduced input flows).
\end{tipbox}
